\section{THREE PROGRAMS, ONE FINDING HISTORY}\label{sec:worlds}

Consider three programs. In each, the finding record shows the pattern
auditors recognize: a program with no finding last year has
about a \nWorldQzero{} chance of a recorded finding this year, and a program
with a finding last year has about a \nWorldQone{} chance. In the notation of
equation~\eqref{eq:qrd}, $q_0 = \nWorldQzero{}$ and $q_1 = \nWorldQone{}$.
Recorded findings are four times as likely after a finding year as after a
clean year. Any audit shop would read that as a strong predictor of another
recorded finding, and would reasonably plan accordingly.

Underneath, the three programs have nothing in common. Table~\ref{tab:worlds}
sets them side by side.

In the first, the condition really did persist. The chance the deficiency is
present rises from $r_0 = 0.100$ after a clean year to $r_1 = 0.400$ after a
finding year, and the audit catches half of what is there in both cases,
$d_0 = d_1 = 0.500$. Multiply: 0.100 times 0.500 is \nWorldQzero{}, and 0.400
times 0.500 is \nWorldQone{}. Equal detection halves the absolute rates and
their difference (a recorded gap of 0.15 against an underlying gap of 0.30)
and leaves the risk ratio exactly as it is underneath.

In the second, nothing about the condition changed at all. The deficiency is
present $r_0 = r_1 = 0.200$ of the time regardless of last year. What changed
is the looking. After a clean year the auditor catches $d_0 = \nWorldOdZero{}$
of what is there; after a finding year, under a follow-up requirement, the
auditor catches all of it, $d_1 = \nWorldOdOne{}$. Multiply: 0.200 times
\nWorldOdZero{} is \nWorldQzero{}, and 0.200 times \nWorldOdOne{} is
\nWorldQone{}. Identical record. The entire fourfold elevation is the
follow-up, and the prior finding carries no information about whether the
condition is there this year.

The third case is the one worth sitting with. Here remediation is stipulated
to halve the underlying rate: the chance the deficiency is present
\emph{falls} from $r_0 = 0.400$ after a clean year to $r_1 = 0.200$ after a
finding year. But the follow-up is thorough. Detection went
from $d_0 = 0.125$ to $d_1 = \nWorldOdOne{}$, a factor of \nWorldRdetRatio{}.
Multiply again and you get \nWorldQzero{} and \nWorldQone{}, the same two
numbers. In the record, programs coming off a finding year look four times
worse than programs coming off a clean year, and in this third case it is the
finding-year programs whose underlying rate is half as high. This column is
a constructed world compatible with successful remediation. It is not evidence
that remediation caused any decline; it shows that the record cannot rule the
possibility out.

The arithmetic is exact, and the three cases are three points on a continuum
along which the observed record is constant.

\subsection*{The observed ratio is a product}

The comparison auditors actually make is a ratio: how much more often findings
are recorded after a prior finding than after a clean year. Three ratios are
needed to say what that comparison contains.

$\Robs = q_1/q_0$ is the \emph{observed finding risk ratio}: how much more
often findings are recorded after a prior finding. In Table~\ref{tab:worlds}
it is 4 in every column, and it is the only one of the three the record
reports.

$\Rlat = r_1/r_0$ is the \emph{underlying-condition risk ratio}: how much
more often the problem itself exists after a prior finding. It is what a
planner would like to know.

$M = d_1/d_0$ is the \emph{history-conditioned detection multiplier}: how
much more likely an existing problem is to be detected in the prior-finding
group than in the clean-history group. Follow-up is the audit mechanism that
motivates it, and the paper calls it the follow-up detection multiplier for
short. It is a ratio of two detection probabilities, not a measure of hours,
sample size or effort.

All three ratios compare groups defined by prior recorded history. None of
them is a before-and-after effect of remediation or of follow-up on a given
program.

Dividing $q_1 = r_1 d_1$ by $q_0 = r_0 d_0$ gives
%
\begin{equation}\label{eq:product}
  \Robs \;=\; \Rlat \,\times\, M .
\end{equation}
%
The ratio of recorded finding rates is the product of the ratio of
underlying-condition rates and the ratio of detection probabilities. The
fourfold finding ratio in Table~\ref{tab:worlds} can come from a fourfold
ratio in underlying problems, a fourfold ratio in detection, or any
combination whose product is four. The table's last three
rows show the three cases: in Persistence, $\Rlat = 4$ and $M = 1$; in
Follow-up only, $\Rlat = 1$ and $M = 4$; in Remediation, $\Rlat = 0.5$ and
$M = 8$. Each multiplies to $\Robs = 4$. A fourfold finding ratio does not by
itself establish fourfold underlying risk, and it does not establish that
follow-up is working either. More years of the same kind of data do not help,
because every additional year is generated by the same product.

\subsection*{Successful remediation can raise recorded findings}

One corollary has a direct operational consequence. Suppose the underlying
condition is \emph{less} common in the prior-finding group, $r_1 < r_0$, as
successful remediation would make it. Recorded
findings still rise after a finding whenever
%
\begin{equation}\label{eq:reversal}
  M \;>\; r_0 / r_1 ,
\end{equation}
%
that is, whenever the follow-up detection multiplier exceeds the factor by
which the problem was reduced. If remediation halves the true problem rate but
follow-up makes detection more than twice as likely, the finding record rises
even though the underlying condition improves. In the Remediation column,
$r_0/r_1 = 2$ and $M = \nWorldRdetRatio{}$, and eight exceeds two. An
organization that judges its corrective-action process by comparing recorded
findings across history groups is using a comparison that can run the wrong
way precisely when the corrective action works and the follow-up is diligent
(Proposition~\ref{prop:remed}).

\begin{table}[tbp]
\centering
\caption{Three programs that produce exactly the same finding record.}
\label{tab:worlds}
\small
\begin{tabular}{lccc}
\toprule
& Persistence & Follow-up only & Remediation \\
\midrule
\multicolumn{4}{l}{\emph{Chance the condition is there, $r$\dots}} \\
\quad after a clean year, $r_0$ & 0.100 & 0.200 & 0.400 \\
\quad after a finding year, $r_1$ & 0.400 & 0.200 & 0.200 \\
\addlinespace
\multicolumn{4}{l}{\emph{Chance the audit detects it if it is there, $d$\dots}} \\
\quad after a clean year, $d_0$ & 0.500 & 0.250 & 0.125 \\
\quad after a finding year, $d_1$ & 0.500 & 1.000 & 1.000 \\
\addlinespace
\multicolumn{4}{l}{\emph{What the record shows, $q = r \times d$\dots}} \\
\quad after a clean year, $q_0$ & 0.05 & 0.05 & 0.05 \\
\quad after a finding year, $q_1$ & 0.20 & 0.20 & 0.20 \\
\addlinespace
\textbf{Observed finding risk ratio}, $R_{\mathrm{obs}} = q_1/q_0$ & \textbf{4.0} & \textbf{4.0} & \textbf{4.0} \\
\textbf{Underlying-condition risk ratio}, $R_{\mathrm{lat}} = r_1/r_0$ & \textbf{4.0} & \textbf{1.0} & \textbf{0.5} \\
\textbf{Detection multiplier}, $M = d_1/d_0$ & \textbf{1.0} & \textbf{4.0} & \textbf{8.0} \\
\bottomrule
\end{tabular}
\tablenotes{Read down each column. In all three, a program with no finding last year has a
5~percent chance of a recorded finding this year and a program with a finding
last year has a 20~percent chance, so the observed finding risk ratio
$R_{\mathrm{obs}}$ is 4 in every column. Underneath, the first program's
condition persisted, the second program's condition was no more likely than
before and only the looking changed, and in the third the underlying rate is
stipulated to be \emph{half} as high after a finding year while detection is
eight times as high. In each column the observed ratio is the product of the
underlying-condition ratio $R_{\mathrm{lat}}$ and the history-conditioned
detection multiplier $M$: $4 = 4 \times 1$,
$4 = 1 \times 4$, $4 = 0.5 \times 8$. Exact arithmetic, not simulation.}
\end{table}

